\documentclass[11pt]{article}
\usepackage[a4paper,margin=27mm]{geometry}
\usepackage{amsmath,amssymb,amsthm,booktabs}
\usepackage[T1]{fontenc}
\usepackage{lmodern}
\usepackage{hyperref}
\newtheorem{proposition}{Proposition}
\title{A three-point finite-field Kakeya set\\\large Counterexample to TLMC 00000009118}
\author{ziangni-sys}
\date{4 October 2026}
\begin{document}
\maketitle
\begin{abstract}
The three nonzero points in $\mathbb F_2^2$ contain an affine line in
every nonzero direction. Thus they form a finite-field Kakeya set of
cardinality $3<2^2$, contradicting the stated lower bound $q^n$.
\end{abstract}

\section{The exact assertion}
The conjecture states that the size of finite-field Kakeya sets has the
explicit lower bound $q^n$. Here $q$ denotes the size of the field and
$n$ the ambient dimension. Neither language version supplies a
multiplicative constant or excludes fields of small cardinality.
We address the asserted coefficient-one bound
\[
 |K|\geq q^n.
\]
This is distinct from a bound $|K|\geq c_n q^n$ with a dimension-dependent
constant $c_n<1$; the example below does not contradict such bounds.

A subset $K\subseteq\mathbb F_q^n$ is a Kakeya set if, for every
nonzero vector $v\in\mathbb F_q^n$, there exists $a\in\mathbb F_q^n$
such that the entire affine line
\[
 \{a+tv:t\in\mathbb F_q\}
\]
lies in $K$. For nonzero $v$, this line has exactly $q$ points.

\section{Explicit counterexample}
Let $q=2$, $n=2$, and work with the field
$\mathbb F_2=\{0,1\}$, with arithmetic modulo two. Define
\[
 K=\mathbb F_2^2\setminus\{(0,0)\}
  =\{(1,0),(0,1),(1,1)\}.
\]
\begin{proposition}
$K$ is a Kakeya set and $|K|=3<q^n=4$.
\end{proposition}
\begin{proof}
The following table lists all nonzero directions and a base point in
each direction. Because the scalar parameter takes only the values
$0$ and $1$, the two rightmost entries give the complete line.
\begin{center}
\begin{tabular}{cccc}
\toprule
Direction $v$ & Base point $a$ & $a+0v$ & $a+1v$\\
\midrule
$(1,0)$ & $(0,1)$ & $(0,1)$ & $(1,1)$\\
$(0,1)$ & $(1,0)$ & $(1,0)$ & $(1,1)$\\
$(1,1)$ & $(1,0)$ & $(1,0)$ & $(0,1)$\\
\bottomrule
\end{tabular}
\end{center}
Every listed point belongs to $K$, and all nonzero directions have
been covered. The three points of $K$ are distinct. Hence $K$ is a
Kakeya set of size three, whereas $q^n=2^2=4$.
\end{proof}

\section{Formal verification}
\raggedright
The Lean project uses Lean 4.19.0 and Mathlib commit
\texttt{c44e0c8ee63ca166450922a373c7409c5d26b00b}.
The ambient space \texttt{Point} is the actual function space
\texttt{Fin 2 -> ZMod 2}; Mathlib supplies the field structure on
\texttt{ZMod 2} and the coordinatewise vector operations.
The definition \texttt{affinePoint} uses ordinary vector addition and
scalar multiplication $a+t\mathbin{\cdot}v$.

The predicate \texttt{IsKakeya} quantifies over all actual nonzero
vectors, existentially quantifies over all actual base points, and
requires membership for every scalar. The set
\texttt{puncturedPlane} is the finite universe with the zero vector
erased. The theorem \texttt{puncturedPlane\_isKakeya} verifies that
full quantified predicate; it does not assume a line-incidence table.
The additional theorem \texttt{affinePoint\_injective} verifies that
every nonzero direction parameterizes two distinct points.

All finite checks use the kernel-checked \texttt{decide} tactic, with
no native evaluation, numerical approximation, or external computation.
The field size, ambient cardinality, and counterexample cardinality
are checked separately. The theorem
\texttt{conjecture\_00000009118} combines the Kakeya property with
the strict inequality. The theorem \texttt{universal\_bound\_false}
explicitly negates the proposed bound for all Kakeya sets in this
ambient space, which suffices to disprove a universal bound over all
finite fields and dimensions.

From the \texttt{lean} directory, run \texttt{lake build}, then
\texttt{lake env lean Main.lean} for the direct verification and
axiom audit. The audited theorems depend only on the standard logical
axioms \texttt{propext}, \texttt{Classical.choice}, and
\texttt{Quot.sound}. No unproved mathematical assumption is used.
\end{document}
